\section{Reproducibility and Evidence Boundaries}
\label{app:repro}

The manuscript accompanies an independent reproducibility bundle in \texttt{sps\_paper/}; it does not modify the earlier \texttt{paper/} draft.
The bundle distinguishes three operations.
\emph{Export} reads completed experiment artifacts and validates the compact evidence snapshot.
\emph{Render} rebuilds the paper's tables and figures from the versioned snapshot without checkpoints, a GPU, shared storage, or network access.
\emph{Plan} describes historical experiment replay and its external prerequisites; it does not submit a job.
Rebuilding a figure is consequently distinct from rerunning the optimization that produced its observations.
The snapshot is stored in \path{data/manifest.json},
\path{data/cohort.json.gz}, \path{data/runs.json.gz}, and
\path{data/probes.json.gz}; \path{data/derived_metrics.json} records the computed summaries.
From the repository root, offline rendering uses
\texttt{python sps\_paper/reproduce.py render}; the nonexecuting historical replay plan uses
\texttt{python sps\_paper/reproduce.py plan}.
The accompanying README records the remaining commands, and the snapshot records complete source commits and full job identifiers.

\paragraph{Design authority.}
Experiment design lives in the Python registries, rather than this manuscript's descriptive summaries.
The ordinary-memory training and objective factorial are defined by
\path{scripts/jobs/run_jobs_sps.py}, including \texttt{TRAINING\_ARGS},
\texttt{lr\_sweep\_experiment\_configs}, and \texttt{objective\_experiment\_configs}.
Shared-prediction training uses \texttt{PredictEmbeddingTrainer} and
\texttt{experiment\_configs} in \path{scripts/analyze_sps_predict_embedding.py}.
Extended-budget partitions and the independent schedule control use the two registries in
\path{scripts/analyze_sps_predict_shards.py}.
The native model's interleaving, shared rotary positions, visibility rule, and odd-slot readout are implemented in
\path{src/compression_horizon/models/sps_causal_lm.py}.
The scalar objectives are implemented in \path{src/compression_horizon/train/loss.py}.
No paper-rendering command launches training.

\paragraph{Historical snapshots.}
The original shared-prediction runs were submitted from
\sourcecommit{9d20b38a8ab842e2dcb97df21cff373c8b99f54a}.
The stretched-schedule shards used
\sourcecommit{f96ca50fe8e834c9fe6f9e0251a15042aaf8a3aa},
and the fixed-trajectory control used
\sourcecommit{cce451fd6625609f1d089de25bb637f479ec827c}.
The objective factorial used
\sourcecommit{4e5726aa444e07b49512cb2af6416a7136b78f3c}.
The route and margin-geometry diagnostics used
\sourcecommit{c155f7c67bda8fbd63d13319322d6f7e4cb2a2e9},
and the local rescue used
\sourcecommit{157a04b98608905364dda5cd311a6c90ab7a3b21}.
These identify historical training/probe code; the current paper-rendering source has its own repository commit.
Complete per-run provenance is retained in the bundle rather than inferred from a run-name prefix.

\paragraph{External prerequisites for replay.}
The Standard and SPS weights are native checkpoints external to the compact paper bundle.
Full replay requires the original environment and GPU implementation.
The plan reconstructs the exact token cache from bundled IDs, trains the upstream ordinary-memory baselines, and links their fresh outputs under the historical diagnostic registry names.
Rescue precedes its dependent gradient audit; checkpoints and cache paths are explicitly relocated.
The texts are the 48 source-cache items selected by the launch registry, not a newly sampled evaluation set.
The snapshot preserves the numerical evidence needed for the reported plots and matched statistics; it is not a substitute for model weights or an independently trained replication.

\paragraph{Integrity checks.}
Selected horizons are matched to the deepest converged raw stage, and target IDs are checked against the pinned cache.
For shared prediction, saved artifacts label the vector's role, initialization, fixed BOS, prediction token, and trainable scalar count.
Generation measurements are recomputed from saved generated and target token IDs.
Shard merges require each registered sample exactly once, complete stage progression, consistent settings and source provenance, and the full target sequence.
For the fixed-trajectory comparison, saved vectors, cumulative counters, and convergence agree exactly with the original control for all 6,342 stages ending before its budget boundary.
These checks concern consistency and provenance of completed measurements, not independent repetition over seeds.

\section{Pretrained Follow-up Evidence and Replay}
\label{app:pretrained}

The additional design registry is \path{scripts/pretrained_sps_experiments.py};
\path{scripts/pretrained_sps_worker.py} implements native-backbone shared-prediction and original-AR training.
The adapter is \path{src/compression_horizon/models/pretrained_sps.py}, with the pinned upstream implementation vendored beside it.
Real-checkpoint preflights check all three variants in each family.
Full-cohort export first invokes the original raw-stage auditor, \path{scripts/analyze_pretrained_sps.py}, before writing \path{data/pretrained.json.gz} and \path{data/pretrained_manifest.json}.
The snapshot contains raw-document hashes/texts, both native token caches with attention masks, selected teacher/greedy token measurements, scalar stage records, and complete job/source provenance.
Optimized vectors and checkpoint weights remain external.

\path{sps_paper/pretrained.py} recomputes the compact snapshot's stage/selection consistency, target identities, stopping reasons and paired statistics offline.
Its \texttt{render} command generates the pretrained tables, figures and \path{data/pretrained_metrics.json}; the standard article build includes this step.
Its \texttt{plan} command prints an executable serial replay without submitting jobs.
The replay restores the exact bundled native caches rather than retokenizing, archives the original worker commits from Git, runs their family preflights, and invokes the original \texttt{train\_shard} with filesystem paths relocated to fresh output and external checkpoint roots.
Training settings are imported from each archived registry rather than retyped into the replay script.
Replay requires original weights, pinned HF snapshots, suitable GPU resources and the archived dependencies; rebuilding the publication assets requires none of them.
CPU checks cover cache restoration, dispatch and offline reconstruction, not a second full GPU optimization of all six cohorts.

\begin{table*}[t]
 \centering
 \small
 % Generated by sps_paper/reproduce.py; edit measurements or generator.
\begin{tabular}{lrr}
\toprule
Paired comparison & $\Delta H$ [95\% CI] & $\Delta G$ [95\% CI] \\
\midrule
Qwen3-1.7B-Base / SPS-CE minus AR & -15.04 [-29.31, +1.25] & -61.71 [-106.60, -16.42] \\
Qwen3-1.7B-Base / SPS-KD minus AR & -19.96 [-33.02, -4.33] & +3.08 [-35.08, +41.21] \\
SmolLM2-1.7B / SPS-CE minus AR & +73.79 [+64.98, +81.60] & -5.35 [-47.90, +38.38] \\
SmolLM2-1.7B / SPS-KD minus AR & +81.31 [+73.25, +89.58] & +33.46 [-9.75, +76.27] \\
\bottomrule
\end{tabular}

 \caption{Pretrained paired differences on the same raw texts within each native-tokenizer family. Percentile 95\% intervals resample paired texts at fixed checkpoints and optimization seed. They do not represent independent training-seed uncertainty or matched-length greedy reliability.}
 \label{tab:pretrained-paired}
\end{table*}

\section{Metrics, Pairing, and Arithmetic}
\label{app:metrics}

The teacher-forced horizon is an integer per text.
It is taken from saved convergence flags and validated by exact target-aligned replay of the selected vector.
Means in the paper are equal-text means, not pooled token-weighted averages.
The initial correct greedy prefix is
\begin{equation}
 G=\max\{g\in\{0,\ldots,H\}:\widehat x_{1:g}=x_{1:g}\},
\end{equation}
where $\widehat x$ is generated using the one selected final vector.
Exact greedy success is $\mathbf{1}\{G=H\}$.
It is not the longest greedily feasible prefix over all saved vectors.
Because different arms select different $H$, their exact greedy counts cannot be interpreted as matched-length reliability measurements.

Paired differences retain the same text ID across compared arms.
The percentile bootstrap resamples those paired IDs with replacement using a fixed analysis seed.
It does not resample model checkpoints, independent optimizations, or pretraining data.
Accordingly, the intervals cannot support population-wide architecture or optimization-seed claims.
We report observed native/random results without choosing an initialization winner from an interval that includes zero.

All main optimizations and free-generation checks use BF16.
Margin deficits are calculated in FP32 after model logits, as implemented in the loss function.
The local geometry and route diagnostics promote the original BF16-quantized weights to FP32 arithmetic; this does not restore information lost in weight quantization.
Their unmodified outputs can differ from the original BF16 training outputs, especially at small margins.
The rescue therefore validates selected vectors again using the actual BF16 forward and free generation.
Even with batch size one, full-prefix teacher replay and incremental generation can take different numerical paths because their sequence shapes differ.
Only directly measured token counts are used for exact reconstruction claims.

\section{Full Ordinary-Memory Objective Factorial}
\label{app:objectives}

\begin{table*}[t]
 \centering
 \small
 % Generated by sps_paper/reproduce.py; edit measurements or generator.
\begin{tabular}{lrrr}
\toprule
Objective / stop & Standard TF [95\% CI] & SPS TF [95\% CI] & Paired gap [95\% CI] \\
\midrule
CE / argmax & 32.31 [30.90, 33.73] & 19.42 [18.29, 20.56] & 12.90 [11.31, 14.58] \\
CE / margin & 32.46 [31.10, 33.81] & 19.44 [18.23, 20.62] & 13.02 [11.60, 14.54] \\
Hinge / argmax & 57.46 [54.85, 60.23] & 35.19 [33.52, 36.88] & 22.27 [19.25, 25.35] \\
Hinge / margin & 54.83 [52.06, 57.73] & 31.52 [29.94, 33.15] & 23.31 [20.10, 26.58] \\
\bottomrule
\end{tabular}

 \caption{Ordinary-memory factorial at the matched original budget. Loss and stopping margin are crossed rather than conflated. The scalar CE and squared-margin objectives are defined in \cref{eq:objective}; raw stopping uses argmax matching and strict stopping requires a positive target gap.}
 \label{tab:factorial}
\end{table*}

The factorial separates two explanations: whether the scalar objective favors every target ranking first, and whether the convergence criterion supplies numerical slack.
Changing raw to strict stopping under CE changes mean lengths only slightly (32.31 to 32.46 for Standard and 19.42 to 19.44 for SPS), while the hinge changes lengths substantially.
The relation need not be monotonic in any single run: changing stopping can alter when a stage ends, its inherited optimizer state, and subsequent updates.
We interpret the complete registered protocol, not a pointwise dominance relation between objectives.

The gradient audit uses four selected old-protocol failed prefixes and saved restart solutions.
At the selected CE vectors, the cosine between CE gradients from currently correct and wrong tokens ranges from $-0.655$ to $-0.560$.
The total gradient norm divided by the sum of the two group norms averages 0.522.
Cosines between negative CE gradient and worst-token margin gradient range from $-0.228$ to $-0.118$; corresponding hinge-descent cosines are 0.868 to 1.000.
These are local derivative measurements of scalar objectives.
They do not establish that actual Adam steps follow plain gradient descent or that SPS has greater gradient conflict than Standard.
The constrained controller's exact solutions have mean CE 1.379, showing that the larger CE of the hinge solutions is not necessary for feasibility.

An earlier CE-only learning-rate grid also leaves a substantial ordinary-memory gap.
\begin{table}[htbp]
 \centering
 \small
 % Generated by sps_paper/reproduce.py; edit measurements or generator.
\begin{tabular}{lrr}
\toprule
Learning rate & Standard TF & SPS TF \\
\midrule
0.01 & 28.85 & 11.77 \\
0.03 & 29.75 & 14.10 \\
0.3 & 32.31 & 19.42 \\
1 & 32.77 & 19.29 \\
\bottomrule
\end{tabular}

 \caption{The four registered CE learning-rate sweep variants. Values are equal-text teacher-forced horizons on the matched cohort; these are measured configurations rather than a globally optimal-rate search.}
 \label{tab:lr-sweep}
\end{table}
Across tested rates, the best observed Standard mean is 32.77 and the best SPS mean is 19.42.
The SPS results at the two largest tested rates differ by only 0.125 tokens with a paired interval including zero.
This motivated a common rate for subsequent controlled comparisons; it does not establish a globally optimal rate or remove selection uncertainty.
The registry and versioned evidence retain the complete grid.

\section{Budget and Terminal-Constraint Audit}
\label{app:terminal}

\begin{table}[htbp]
 \centering
 \small
 \resizebox{\columnwidth}{!}{% Generated by sps_paper/reproduce.py; edit measurements or generator.
\begin{tabular}{lr}
\toprule
Final-stage measurement & Value \\
\midrule
Median margin violations & 10 \\
Range of margin violations & 2--92 \\
Mean fraction meeting margin & 0.871 \\
Previous stage converged & 48/48 \\
Minimum old-position violations & 1 \\
Median old-position lower bound & 9 \\
\bottomrule
\end{tabular}
}
 \caption{Saved terminal stages of the native shared-prediction fixed-trajectory extension. Violations concern the recorded positive-margin convergence criterion. The old-constraint count is a lower bound obtained by allowing the single added target to account for one violation.}
 \label{tab:terminal}
\end{table}

If a failed stage extends $L-1$ converged targets to $L$ and records $V$ violations of the positive-margin criterion, at most one can be the newly added target.
The number of old violations is therefore at least $V-1$.
The observed minimum $V=2$ establishes at least one old violation in every text; the median $V=10$ establishes a median lower bound of nine.
This reasoning uses the saved counts and stage relationship only.
It neither identifies individual positions nor says that all violated margins correspond to incorrect argmax predictions.

The final saved vector has changed during the failed optimization, and extending the context can also change BF16 computation.
Without per-position evaluation at the warm start and optimized vector under controlled shapes and arithmetic, those explanations cannot be separated.
The audit therefore supports simultaneous constraint loss, while leaving the cause and recoverability at the new frontier unresolved.

Saved-stage optimizer costs also help interpret the original native-hinge shared-prediction run.
Median cumulative updates at positions 10, 25, 50, 75, 100 and 125 are 179, 529.5, 1,269, 3,316.5, 6,344 and 8,476.
The contributing cohorts are 48, 48, 48, 48, 47 and 38 texts.
Median incremental successful-stage costs at those positions are 17, 33.5, 40, 70, 52 and 71 updates.
They are noisy rather than monotonic and are costs of reoptimizing the whole prefix with inherited optimizer state and a running schedule, not isolated-token complexity.
Zero-update successes are retained.
Late summaries condition on surviving texts and cannot establish length as the sole causal driver of cost.

\section{Ordinary-Memory Intervention and Rescue Details}
\label{app:mechanism}

\begin{table}[htbp]
 \centering
 \small
 \resizebox{\columnwidth}{!}{% Generated by sps_paper/reproduce.py; edit measurements or generator.
\begin{tabular}{lrr}
\toprule
Intervention & Standard $\Delta$NLL & SPS $\Delta$NLL \\
\midrule
All memory reads removed & 4.546 & 4.456 \\
Context reads removed & -- & 3.273 \\
Predict reads removed & -- & 4.221 \\
Donor keys, all layers & 2.396 & 2.450 \\
Donor values, all layers & 5.277 & 4.834 \\
Donor keys + values & 5.220 & 4.944 \\
Self patch, keys + values & 0.000 & 0.000 \\
\bottomrule
\end{tabular}
}
 \caption{Ordinary-memory causal route and donor-patch summaries. Changes are equal-text mean NLL relative to each model's own baseline on targets after the first. Layer indices are one-based. Key and value patches are counterfactual interventions rather than additive information allocations.}
 \label{tab:routes}
\end{table}

\begin{table}[htbp]
 \centering
 \small
 % Generated by sps_paper/reproduce.py; edit measurements or generator.
\begin{tabular}{lrr}
\toprule
Restart & FP32 exact & BF16 exact \\
\midrule
CE Adam (0.3) & 0/4 & 0/4 \\
CE Adam (1.0) & 0/4 & 0/4 \\
Hinge Adam & 4/4 & 3/4 \\
Trust-region margin & 4/4 & 1/4 \\
\bottomrule
\end{tabular}

 \caption{Local restart results at four selected failed SPS ordinary-memory prefixes, without added trainable scalars. FP32 uses the original BF16-quantized weights promoted for arithmetic. Original-regime success requires actual BF16 teacher replay and greedy reconstruction. These are prefix rescues, not full progressive runs.}
 \label{tab:rescue}
\end{table}

The route probe uses 16 texts at target length 12 and memories optimized for exactly that stage, rather than cropping longer-stage vectors.
Donor memories come from another selected text and are rescaled to match the recipient norm before computing their projected keys/values.
The recipient residual trajectory is preserved during a key-only or value-only patch.
Removing memory-key edges changes softmax normalization and later states, so joint and individual effects need not add.
The broader attention-ablation probe likewise measures sensitivity of already optimized memories, not a derivative of achievable horizon with respect to a route.

The local geometry probe uses eight selected texts at their first failed SPS ordinary-memory length, evaluating both checkpoints at that same length and at warm-start and optimized memories.
For the old-token margin Jacobian $J_{<L}$ and new-token gradient $g_L$, it measures
\begin{equation}
 \rho=\frac{\|P_{\ker J_{<L}}g_L\|_2}{\|g_L\|_2}.
\end{equation}
Full old-row rank and a substantial $\rho$ leave a first-order direction that changes the new margin without changing old margins.
They provide no guarantee of a finite nonlinear feasible move.
The median minimum previous-token margins at warm start are 0.147 for Standard and 0.029 for SPS.
Because the sampling frontier differs between models, this smaller SPS slack may be a consequence of proximity to its stopping point.
It is not an isolated architecture-specific cause.
Likewise, PCA of a progressive optimization trajectory describes the path taken by the optimizer rather than the dimension of the feasible solution set; we do not use trajectory dimensionality as a capacity certificate.

The four-text rescue compares CE Adam restarts, squared-margin Adam, and a constrained margin controller from saved failed-stage memories.
The controller solves a local minimum-norm margin problem with a trust region and finite-forward validation.
Diagnostic accounting treats one scalar-loss backward as one vector--Jacobian product and a full $L$-row Jacobian as $L$ such products; forward counts are tracked separately.
Equal derivative budgets are not equal wall-clock budgets.
Squared-margin Adam reaches all four FP32 solutions in 20--55 gradients, but only three retain both exact criteria in BF16.
No result here is a rescue attempt at the final shared-prediction frontier.

\section{Numerical stability across prefix lengths and cached decoding}
\label{app:numerics}

Causal masking makes the next-token distribution a function of the preceding tokens.
In floating-point arithmetic, evaluating those tokens in a longer tensor can change the order of a reduction.
A change in the winning logit at one position is sufficient for blind greedy replay to depart from the target.
We therefore measure two quantities separately: the number of unequal vocabulary logits on identical token histories, and the number of positions at which the highest-scoring token changes.
All comparisons below use batch size one. Teacher-forced comparisons always feed the same tokens to both paths.

\paragraph{Prefix length and cache boundaries.}
We evaluated the frozen BF16 Qwen3-1.7B-Base model on a single PG-19 passage of 32,768 tokens using PyTorch SDPA.
For a prefix of length $L$, we compared its first $L$ output positions with those same positions in the longer pass.
At $L=2{,}048$, the winner changed at 44 positions. At $L=8{,}192$ the shared logits agreed throughout, yet splitting that same 8,192-token input into a 64-token prefill and a cached suffix changed 159 winners relative to a single pass (\cref{tab:numerics-length}).
Thus agreement between two uncached lengths did not predict agreement across a cache boundary on the same text.

\begin{table*}[htbp]
 \centering\small
 \begin{tabular}{lrrrrrrrr}
\toprule
Length & Flips & MAE & Top-3 & Top-5 & Top-10 & Step & Split & Split vs. 32k \\
\midrule
128 & 2 & 0.0405 & 97.40 & 98.44 & 98.75 & 1 & 2 & 2 \\
512 & 10 & 0.0450 & 97.92 & 97.58 & 98.20 & 0 & 5 & 6 \\
1,024 & 20 & 0.0449 & 97.59 & 97.87 & 98.32 & 1 & 14 & 18 \\
2,048 & 44 & 0.0425 & 97.77 & 98.15 & 98.37 & 1 & 36 & 40 \\
4,096 & 0 & 0.0000 & 100.00 & 100.00 & 100.00 & 0 & 75 & 0 \\
8,192 & 0 & 0.0000 & 100.00 & 100.00 & 100.00 & 0 & 159 & 0 \\
32,768 & 0 & 0.0000 & 100.00 & 100.00 & 100.00 & 1 & 665 & 0 \\
\bottomrule
\end{tabular}

 \caption{Prefix and cache effects on Qwen with default SDPA. The first five metrics compare an uncached prefix with the corresponding rows of an uncached 32,768-token pass. Flips count changed winners, MAE averages absolute logit differences over positions and vocabulary, and top-$k$ reports mean set overlap as a percentage of $k$ (token ID breaks ties). The last three columns count changed winners. Step compares the last 32 positions decoded one at a time with the same-length uncached pass. Split compares all $L-64$ suffix positions following a 64-token prefill with the same-length uncached pass. Split vs.\ 32k compares the shared suffix rows of two cached passes with the same split point. The 32,768-token row includes the reference self-comparisons.}
 \label{tab:numerics-length}
\end{table*}

\paragraph{Batch-invariant operations and attention.}
On the first 8,192 tokens of the same passage, we enabled the \href{https://github.com/thinking-machines-lab/batch_invariant_ops/tree/f22b1fbe534b1fe045080f44a970a441caaf3c4f}{Thinking Machines batch-invariant PyTorch operations} \citep{he2025batchinvariance} in every arm and varied the attention implementation.
For the split after 64 tokens, default SDPA changed 159 suffix winners with the stock operators and 144 with the batch-invariant operators. In the latter run, the first layer's shared queries, keys and values agreed exactly. Default SDPA changed 1,208,280 elements of that layer's suffix attention output.
Our Triton forward anchors query tiles to absolute positions and visits key blocks in a fixed order across prefill and decode.
It reproduced its own full-pass logits at each tested split and in the one-token test (\cref{tab:numerics-attention}).
Changing the kernel also changed the full-pass reference: fixed tiles and default SDPA selected different winners at 145 of 8,192 positions.
All three arms ran on an A100 using PyTorch 2.8.0 and Triton 3.4.0.

\begin{table*}[htbp]
 \centering\small
 \begin{tabular}{lrrrr}
\toprule
Attention & Split at 64 & Split at 65 & Split at 8,191 & Last 32 steps \\
\midrule
Default SDPA & 879,297,654 / 144 & 877,954,503 / 143 & 111,711 / 0 & 3,602,909 / 0 \\
Math SDPA & 873,519,051 / 147 & 873,417,988 / 147 & 100,693 / 0 & 3,581,841 / 0 \\
Fixed tiles & 0 / 0 & 0 / 0 & 0 / 0 & 0 / 0 \\
\bottomrule
\end{tabular}

 \caption{Unequal vocabulary logits / positions with a changed top-1 token. Each cell compares cached outputs with the same backend's full pass. The split columns score 8,128, 8,127 and one suffix position, respectively. The final column scores 32 successive single-token forwards on the true history. Each row includes the Thinking Machines operators.}
 \label{tab:numerics-attention}
\end{table*}

\paragraph{Native SPS cramming.}
The combined training backend used the Thinking Machines operators, fixed-tile attention with an FP32 row sum, and an analytical attention backward pass. Relative to the stock backend, exact blind recovery increased from 38/48 to 45/48 texts for SPS and from 44/48 to 48/48 for Standard, at each run's own selected horizon.
For the three remaining SPS failures, a change in mask-stride alignment changed how Triton grouped the FP32 additions in the softmax denominator.
At the first affected layer, the attention numerators agreed on every shared element.
Widening just the per-block sum to FP64 and casting it back to FP32 recovered these three archived vectors.
Across all 48 archived SPS vectors, this change also moved ten other vectors across a greedy decision boundary in the full teacher-forced pass.
With the wider sum, full, growing and cached evaluation agreed on all 315,657,600 logits at 6,275 target positions. All 38 vectors with exact teacher-forced replay under that arithmetic recovered their targets in blind cached decoding.
The weights remained BF16 throughout.
On the original and fixed-tile FP32-sum training runs, mean wall time per logged optimization step rose from 29.76 to 47.74 ms for Standard and from 29.79 to 49.75 ms for SPS. These timings include the training loop and the attention backward pass.
